% ==========================================================================
% GENERATED by scripts/split_paper.py. Do not edit this file.
%
% Source: imported_papers/Neural_Computation__CNA/NECO-D-24-00304-Source.tex
% Part 12: Single Weight-Kernel Implementation
%
% Edits belong upstream, in the published paper. To change what the
% chapter includes, change imports/cna.toml and re-run
% the splitter.
% ==========================================================================
\section{Single Weight-Kernel Implementation}
So far, we introduced the weights $\boldsymbol{W}^{(F)}$ and $\boldsymbol{W}^{(L)}$ in $S2$ as two independent convolutional kernels.
For simplicity in our algorithm, and since we use the same kernel size for the forward and recurrent connections, we stack the two kernels as $\boldsymbol{W}^{(S2)} = (\boldsymbol{W}^{(F)}; \boldsymbol{W}^{(L)})$ as well as the inputs $\boldsymbol{y}^{(S1)}$ and $\boldsymbol{y}^{(S2)}_{t-1}$, leading to the equivalent formulation of eq. \ref{eq:a_a2}:

\begin{equation}
  \boldsymbol{a}^{\mathcal{I}(S2)}_{c_{\text{out}},t} = \sum_{c_\text{in}=0}^{C^{(S2)}} \boldsymbol{W}^{(S2)}_{c_{\text{out}}, c_{\text{in}}} \star \left( \boldsymbol{y}_{c_{\text{in}}}^{(S1)}; \boldsymbol{y}_{c_{\text{in}},t-1}^{(S2)} \right)
  \label{eq:a_a22}
\end{equation}
